\documentclass[14pt]{extarticle}
% ---------- PACKAGES ----------
\usepackage{amsmath, amssymb}
\usepackage{geometry}
\usepackage{setspace}
\usepackage{titlesec}
\usepackage{xcolor}
\usepackage{helvet}
\usepackage{enumitem}
\usepackage{lmodern}
\makeatletter
\IfFileExists{../common-things/reflectionpage.tex}{%
  \input{../common-things/reflectionpage.tex}%
}{%
  \IfFileExists{common-things/reflectionpage.tex}{%
    \input{common-things/reflectionpage.tex}%
  }{%
    \PackageError{\jobname}{Missing reflectionpage.tex file}{%
      Place reflectionpage.tex inside a common-things/ directory next to this unit\@.%
    }%
  }%
}
\makeatother
% ---------- PAGE SETUP ----------
\geometry{margin=1in}
\setstretch{1.5}
\setlength{\parindent}{0pt}
\setlength{\parskip}{0.75em}
\setlist{itemsep=0.75\baselineskip, topsep=0.5\baselineskip}
\titleformat{\section}{\normalfont\Large\bfseries}{\thesection}{1em}{}
\titleformat{\subsection}{\normalfont\large\bfseries}{\thesubsection}{1em}{}
\renewcommand{\familydefault}{\sfdefault}
\renewcommand{\emph}[1]{\textbf{#1}}

\begin{document}
\raggedright
\pagenumbering{gobble}

\begin{center}
    \LARGE \textbf{Unit 10: Review and Test Strategy} \\[6pt]
    \Large \textbf{Topic 4: Pacing and the Two-Pass Approach}
\end{center}

\vspace{1em}

\section*{Concept Summary}

The Digital SAT Math section consists of two modules of 22 questions each, with 35 minutes per module. That gives you approximately \textbf{1 minute and 35 seconds per question} on average. Some questions take 30 seconds; others take 3 minutes. The key to strong performance is not spending equal time on every question but rather managing your time so that you attempt every question you are capable of answering correctly.

\textbf{The Two-Pass Approach:}

\textbf{Pass 1 (first 20--25 minutes):} Work through every question in order. If a question feels straightforward, solve it and move on. If a question looks time-consuming or confusing after 30 seconds of reading, \textbf{flag it and skip it}. Do not spend 4 minutes on one question while easier questions sit unanswered later in the module. The Digital SAT interface has a built-in flag/bookmark feature for exactly this purpose.

\textbf{Pass 2 (remaining 10--15 minutes):} Return to every flagged question. With the easy points already secured, you can invest more time in the harder problems without anxiety. You may also find that a later question triggered an idea relevant to an earlier one.

\textbf{Pacing benchmarks:} After 15 minutes, you should have completed roughly 10--12 questions. After 25 minutes, you should have at least a first-pass attempt on all 22 questions. If you are significantly behind these marks, you are spending too long on individual problems.

\textbf{Adaptive difficulty:} The Digital SAT is adaptive. If you do well on Module 1, Module 2 will be harder but worth more toward your score. This means Module 2 questions are expected to take longer. Adjust your pacing accordingly: you might need a more aggressive skip threshold on Module 2.

\textbf{Guessing strategy:} There is no penalty for wrong answers on the Digital SAT. Never leave a question blank. If you are running out of time, enter your best guess for every remaining question. For free-response questions, common ``default guesses'' include 0, 1, 2, or other small integers, since SAT answers tend to be simple values.

\section*{Core Skills}
\begin{itemize}
  \item Classify questions into ``quick solve,'' ``medium effort,'' and ``skip for now'' within 15--20 seconds of reading.
  \item Use the flag/bookmark feature systematically during Pass 1.
  \item Monitor your pacing against time benchmarks (roughly half the questions done by the halfway mark).
  \item Allocate remaining time to flagged questions in order of confidence, not in order of appearance.
  \item Never leave a free-response answer blank; enter a reasonable guess.
\end{itemize}

\section*{Example 1: Quick-Solve Recognition}

\textit{``If \(3x = 18\), what is \(x + 5\)?''}

This is a 15-second question. Solve \(x = 6\), then \(x + 5 = 11\). Do not overthink it. Answer and move on.

\(\boxed{11}\)

\section*{Example 2: Medium-Effort Recognition}

\textit{``A rectangle has a perimeter of 30 and an area of 50. What is the length of the longer side?''}

This requires setting up a system (\(l + w = 15\), \(lw = 50\)) and solving a quadratic. It is doable in about 90 seconds. Solve on Pass 1 if you are on pace, or flag it if you are behind.

\textbf{Solution:} \(l + w = 15\) and \(lw = 50\). Substituting \(w = 15 - l\): \(l(15 - l) = 50 \implies l^2 - 15l + 50 = 0 \implies (l - 5)(l - 10) = 0\). The longer side is 10.

\(\boxed{10}\)

\section*{Example 3: Skip-and-Return Recognition}

\textit{``A function \(f\) satisfies \(f(x+2) = f(x) + 3\) for all \(x\). If \(f(1) = 4\), what is \(f(9)\)?''}

This requires careful thought: how many steps of \(+2\) get from 1 to 9? If it is not immediately clear, flag it and return on Pass 2.

\textbf{Solution:} From \(x = 1\) to \(x = 9\) is 4 steps of \(+2\). Each step adds 3 to the output, so \(f(9) = 4 + 4(3) = 16\).

\(\boxed{16}\)

\section*{Key Takeaways}
\begin{itemize}
  \item The Two-Pass strategy ensures you see every question and prioritize the ones you can solve confidently. Spending 5 minutes on one hard question while skipping two easy ones is the most common pacing mistake.
  \item Use the 15-minute checkpoint: if you have not completed at least 10 questions, speed up your skip threshold.
  \item On Pass 2, tackle flagged questions in order of confidence, not in the order they appear. Start with the one you feel best about.
  \item Never leave a question blank. On free-response, guess a small integer if you have no idea. On the SAT, the most common answers are small positive numbers.
\end{itemize}

\newpage

\section*{Practice Questions: Pacing and the Two-Pass Approach}

The questions below are a timed mini-set. Set a timer for \textbf{14 minutes} and attempt all 20 questions. Use the Two-Pass approach: do a quick first pass, flag hard ones, then return. Track which questions you flagged and how long each pass took.

\subsection*{Part A: Quick Solves (target: under 30 seconds each)}
\begin{enumerate}
  \item If \(4x = 28\), what is \(x\)?
  \item What is \(15\%\) of 80?
  \item A triangle has angles measuring \(45^\circ\), \(90^\circ\), and \(x^\circ\). What is \(x\)?
  \item If \(f(x) = 3x - 2\), what is \(f(5)\)?
  \item The mean of 3, 7, and \(x\) is 6. What is \(x\)?
\end{enumerate}

\subsection*{Part B: Medium Effort (target: 60--90 seconds each)}
\begin{enumerate}
  \item Solve the system: \(x + y = 10\) and \(2x - y = 5\).
  \item A circle has an area of \(49\pi\). What is its circumference? Express your answer in terms of \(\pi\).
  \item Simplify: \(\dfrac{x^2 - 4}{x + 2}\) for \(x \neq -2\).
  \item A shirt costs \(\$40\) after a \(20\%\) discount. What was the original price?
  \item The exterior angle of a triangle is \(120^\circ\). One remote interior angle is \(45^\circ\). What is the other remote interior angle, in degrees?
\end{enumerate}

\subsection*{Part C: Challenging (target: 90--120 seconds each)}
\begin{enumerate}
  \item A right triangle has a hypotenuse of 13 and one leg of 5. What is the area of the triangle?
  \item For what value of \(k\) does \(x^2 - 8x + k = 0\) have exactly one solution?
  \item A cylindrical can has a radius of 3 cm and a volume of \(108\pi\) cubic cm. What is the height, in cm?
  \item The function \(f(x) = -2(x - 3)^2 + 18\) has a maximum value. What is it?
  \item If \(\sin(\theta) = \dfrac{5}{13}\) and \(\theta\) is acute, what is \(\cos(\theta)\)?
\end{enumerate}

\subsection*{Part D: Hard (target: 2--3 minutes each)}
\begin{enumerate}
  \item A store marks up an item by \(60\%\) and then discounts it by \(25\%\). What is the net percent change from the original cost?
  \item The line \(y = mx\) is tangent to \(y = x^2 + 1\). Find the positive value of \(m\).
  \item In a 30-60-90 triangle, the side opposite the \(60^\circ\) angle is \(8\sqrt{3}\). What is the perimeter of the triangle?
  \item A weighted average: 20 students scored a mean of 70 and 30 students scored a mean of 85. What is the overall mean?
  \item Two similar triangles have areas in the ratio \(4:9\). If the perimeter of the larger triangle is 45, what is the perimeter of the smaller triangle?
\end{enumerate}

\newpage

\section*{Answer Key and Solutions: Pacing and the Two-Pass Approach}

\subsection*{Part A Solutions: Quick Solves}
\begin{enumerate}
  \item \(x = 7\).

  \(\boxed{7}\)

  \item \(0.15 \times 80 = 12\).

  \(\boxed{12}\)

  \item \(x = 180 - 45 - 90 = 45\).

  \(\boxed{45}\)

  \item \(f(5) = 15 - 2 = 13\).

  \(\boxed{13}\)

  \item \(\dfrac{3 + 7 + x}{3} = 6 \implies 10 + x = 18 \implies x = 8\).

  \(\boxed{8}\)
\end{enumerate}

\subsection*{Part B Solutions: Medium Effort}
\begin{enumerate}
  \item Add the equations: \(3x = 15 \implies x = 5\). Then \(y = 5\).

  \(\boxed{(5, 5)}\)

  \item \(\pi r^2 = 49\pi \implies r = 7\). Circumference \(= 2\pi(7) = 14\pi\).

  \(\boxed{14\pi}\)

  \item \(\dfrac{(x-2)(x+2)}{x+2} = x - 2\).

  \(\boxed{x - 2}\)

  \item \(0.80p = 40 \implies p = 50\).

  \(\boxed{\$50}\)

  \item The other remote interior angle is \(120 - 45 = 75^\circ\).

  \(\boxed{75}\)
\end{enumerate}

\subsection*{Part C Solutions: Challenging}
\begin{enumerate}
  \item Other leg: \(\sqrt{169 - 25} = 12\). Area: \(\dfrac{1}{2}(5)(12) = 30\).

  \(\boxed{30}\)

  \item Discriminant \(= 0\): \(64 - 4k = 0 \implies k = 16\).

  \(\boxed{16}\)

  \item \(\pi(9)(h) = 108\pi \implies 9h = 108 \implies h = 12\).

  \(\boxed{12}\)

  \item The maximum value is the \(k\) in vertex form: 18.

  \(\boxed{18}\)

  \item \(\cos(\theta) = \sqrt{1 - 25/169} = \sqrt{144/169} = \dfrac{12}{13}\).

  \(\boxed{\dfrac{12}{13}}\)
\end{enumerate}

\subsection*{Part D Solutions: Hard}
\begin{enumerate}
  \item Net factor: \(1.60 \times 0.75 = 1.20\). Net change: \(20\%\) increase.

  \(\boxed{20\%}\)

  \item Set \(mx = x^2 + 1\), giving \(x^2 - mx + 1 = 0\). For tangency, discriminant \(= 0\):
  \[
  m^2 - 4 = 0 \implies m = 2 \text{ (taking the positive value)}
  \]

  \(\boxed{2}\)

  \item The side opposite \(60^\circ\) is \(s\sqrt{3} = 8\sqrt{3}\), so \(s = 8\). The shortest side is 8, the side opposite \(60^\circ\) is \(8\sqrt{3}\), and the hypotenuse is \(2(8) = 16\). Perimeter: \(8 + 8\sqrt{3} + 16 = 24 + 8\sqrt{3}\).

  \(\boxed{24 + 8\sqrt{3}}\)

  \item Weighted mean: \(\dfrac{20(70) + 30(85)}{50} = \dfrac{1400 + 2550}{50} = \dfrac{3950}{50} = 79\).

  \(\boxed{79}\)

  \item Side ratio: \(\sqrt{4/9} = 2/3\). Perimeter of smaller: \(45 \times \dfrac{2}{3} = 30\).

  \(\boxed{30}\)
\end{enumerate}

\ReflectionPage{Unit 10, Topic 4}{https://www.scotchegg.co}

\end{document}
